\begin{lemma}[Input-independent public preprocessing]\label{lem:step-010-public-preprocessing}
Under Assumption~\ref{assump:candidate-regime}, the contradiction hypothesis~\eqref{eq:local-negation},
Propositions~\ref{prop:step-004-finite-hard-prior},
\ref{prop:step-005-certificate}, and
\ref{prop:step-009-overflow}, the following ordered procedure defines a
preprocessing seed \(\Omega_{\mathrm{pre}}\) before any input coordinate is
evaluated:

\begin{enumerate}
\item Draw \(J\sim\operatorname{Unif}[k]\).
\item Draw \(I_1,\ldots,I_n\) independently and uniformly from \([k]\),
independently of \(J\).
\item Conditional on \(J\), independently draw
   \[
     \Theta_i=(T_i,Q_i)\sim\mu_{N,M},
     \qquad i\in[k]\setminus\{J\};
   \]
\item Conditional on \(J,I_{1:n},(\Theta_i)_{i\ne J}\), independently draw
   \[
     X_r^\circ\sim Q_{I_r}
     \quad\text{for every }r\in[n]\text{ with }I_r\ne J,
   \]
   and set
   \[
     Y_r^\circ:=\tau_{T_{I_r}}(X_r^\circ).
   \]
\end{enumerate}

The seed consists of precisely these realized objects. Its distribution
depends only on the public tuple
\((k,N,n,M,\mu_{N,M})\), is independent of every
\(z\in\mathsf Z_M\), and contains neither a hidden instance
\((t,Q)\), a block-instance draw \(\Theta_J\), nor a synthetic feature or
label at a position whose tag equals \(J\). Conditional on the realized
selector, tags, and nonhidden instances, the nonhidden labeled
rows

\[
  \left((I_r,X_r^\circ),Y_r^\circ\right),
  \qquad r:I_r\ne J,
\tag{6}
\]

are mutually independent, and each row with tag \(i\ne J\) has labeled law
\(Q_i^{\tau_{T_i}}\) after its tag is suppressed.
\end{lemma}

\begin{proof}
Proposition~\ref{prop:step-005-certificate} discharges the exact
\(N,M\) conditions for Proposition~\ref{prop:step-004-finite-hard-prior}, so the finite public law
\(\mu_{N,M}\) exists before the present simulator is defined. The first two
draws use the selector/tag product kernel certified by Proposition~\ref{prop:step-009-overflow}. Once \(J=j\) is fixed, the finite
product law

\[
  \bigotimes_{i\in[k]\setminus\{j\}}\mu_{N,M}
\]

defines all nonhidden instances. Each realized \(Q_i\) is a probability
distribution on the finite set \([N]\), so the further finite conditional
product

\[
  \bigotimes_{\substack{r\in[n]\\ I_r\ne J}}Q_{I_r}
\tag{7}
\]

defines the nonhidden feature draws. For any fixed values
\((x_r)_{r:I_r\ne J}\), (7) gives the exact conditional mass

\[
  \prod_{\substack{r\in[n]\\ I_r\ne J}}Q_{I_r}(x_r),
\tag{8}
\]

which proves conditional mutual independence and the stated marginal
kernels. Applying the deterministic threshold label
\(\tau_{T_{I_r}}\) gives exactly the labeled law in the lemma.

All four randomization stages are completed without referring to an input
argument. In particular, the sampling code knows the public distribution
\(\mu_{N,M}\) and its own realized nonhidden pairs
\((T_i,Q_i)\), but it neither draws nor receives \((T_J,Q_J)\).
A later hidden pair \((t,Q)\) is therefore not an input to this procedure.
Endpoint draws \(T_i=1\) and \(T_i=N+1\) are executable: their deterministic
labels are respectively all one and all zero. Point-mass \(Q_i\)'s are also
valid factors in (7).

The internal randomness used by \(A\) is intentionally not a component of
\(\Omega_{\mathrm{pre}}\). It remains the fresh internal randomization of
the one call to the kernel \(A\), if that call is reached. Thus every
random choice belonging to preprocessing is fixed before input access,
without converting randomized \(A\) into a deterministic map.
\end{proof}

\begin{proposition}[Exact hidden-blind size-\(n\) row construction]\label{prop:step-010-row-construction}
Under Lemma~\ref{lem:step-010-public-preprocessing}, fix any realization
\(\omega\) of \(\Omega_{\mathrm{pre}}\) and any arbitrary labeled input

\[
  z=(z_1,\ldots,z_M)\in\mathsf Z_M,
  \qquad z_\ell=(x_\ell,y_\ell).
\tag{9}
\]

Define, using only the selector and tags in \(\omega\),

\[
  L_r:=\sum_{s=1}^r\mathbf1\{I_s=J\},
  \qquad
  U:=L_n.
\tag{10}
\]

The value of \(U\) is computed before any \(z_\ell\) is evaluated. If
\(U>M\), the row-construction procedure stops immediately and evaluates no
input coordinate. If \(U\le M\), it defines an ordered size-\(n\) labeled
product dataset

\[
  \mathcal P_\omega(z)
  =(\widetilde z_1,\ldots,\widetilde z_n)\in\mathsf S_n
\tag{11}
\]

by the exact rule

\[
  \widetilde z_r=
  \begin{cases}
    ((J,x_{L_r}),y_{L_r}), & I_r=J,\\[2mm]
    ((I_r,X_r^\circ),Y_r^\circ), & I_r\ne J.
  \end{cases}
\tag{12}
\]

Thus, on nonoverflow, the \(\ell\)-th input row is used exactly at the
\(\ell\)-th hidden-tag occurrence for every \(\ell\in[U]\), every
nonhidden row is the pre-drawn row (6), and the rule requires no knowledge
of a hidden \(t\) or \(Q\). It is valid for every binary label pattern in
(9), including patterns inconsistent with every threshold.
\end{proposition}

\begin{proof}
The seed is drawn before the input is accessed by
Lemma~\ref{lem:step-010-public-preprocessing}. Both \(L_r\) and \(U\) are
functions only of the already drawn \(J,I_1,\ldots,I_n\). Consequently the
comparison \(U>M\) can be resolved before evaluating any coordinate of
(9). On overflow, the procedure takes no later row-construction action, so
it never attempts to access a nonexistent \((M+1)\)-st record.

Suppose \(U\le M\). For each position \(r\) with \(I_r=J\), (10) gives

\[
  1\le L_r\le U\le M.
\tag{13}
\]

Hence \(z_{L_r}\) is an available input record. Since
\(x_{L_r}\in[N]\), \(J\in[k]\), and \(y_{L_r}\in\{0,1\}\), the first line
of (12) lies in \(X_{k,N}\times\{0,1\}\), regardless of whether
\(y_{L_r}=\tau_t(x_{L_r})\) for any \(t\).

For each position with \(I_r\ne J\), the seed already contains
\(X_r^\circ\in[N]\) and
\(Y_r^\circ=\tau_{T_{I_r}}(X_r^\circ)\in\{0,1\}\), so the second line of
(12) has the same required type. There is one and only one applicable line
of (12) for every \(r\in[n]\). Thus (11) has exactly \(n\) ordered labeled
rows and belongs to the full domain \(\mathsf S_n\) of \(A\).

If \(r\) is the \(\ell\)-th occurrence of tag \(J\), then by definition
\(L_r=\ell\), so the first line of (12) uses \(z_\ell\), including its raw
label. This is why no hidden threshold \(t\) is needed. The feature
distribution \(Q\) of a later hidden input law is likewise unnecessary:
the procedure takes hidden features directly from the input records.
Nonhidden features and labels use only the simulator's publicly drawn
instances, whose exact conditional kernels were proved in
Lemma~\ref{lem:step-010-public-preprocessing}.
\end{proof}

\begin{lemma}[Ordered one-use input incidence]\label{lem:step-010-one-use}
Under Proposition~\ref{prop:step-010-row-construction}, for every fixed
preprocessing seed \(\omega\) and every arbitrary input
\(z\in\mathsf Z_M\), the following statements hold pathwise.

\begin{itemize}
\item If \(U>M\), no coordinate of \(z\) is evaluated.
\item If \(U\le M\), let
  \[
    r_\ell:=\min\left\{r\in[n]:
      \sum_{s=1}^r\mathbf1\{I_s=J\}=\ell\right\},
    \qquad \ell=1,\ldots,U.
  \tag{14}
  \]
  Then
  \[
    1\le r_1<\cdots<r_U\le n,
    \qquad
    \widetilde z_{r_\ell}=((J,x_\ell),y_\ell).
  \tag{15}
  \]
  No row of \(\mathcal P_\omega(z)\) other than row \(r_\ell\) depends on
  input coordinate \(z_\ell\); coordinates \(z_{U+1},\ldots,z_M\) are not
  evaluated at all.
\end{itemize}

In particular, each input record has at most one image row. If \(U=0\), the
input-access set is empty and all \(n\) rows are nonhidden pre-drawn rows.
If \(U=M\), every input record has exactly one image row. The statement is
about row incidence only and makes no differential-privacy claim.
\end{lemma}

\begin{proof}
On overflow, the first bullet is part of the control flow proved in
Proposition~\ref{prop:step-010-row-construction}. Now fix \(U\le M\). The
running count \(L_r\) in (10) starts at zero, is nondecreasing, and
increases by exactly one precisely at the positions with \(I_r=J\).
Because its terminal value is \(U\), it hits each integer
\(\ell=1,\ldots,U\) at a unique first position \(r_\ell\), and those
positions are strictly increasing. At \(r_\ell\), one has
\(I_{r_\ell}=J\) and \(L_{r_\ell}=\ell\), so the first line of (12) gives
the identity in (15).

At a hidden position \(r\), the only input coordinate named by (12) is
\(z_{L_r}\). The map from hidden positions to positive running counts is
the bijection \(r_\ell\mapsto\ell\). At a nonhidden position, the second
line of (12) is already fixed by \(\omega\) and names no input coordinate.
Therefore \(z_\ell\), for \(\ell\le U\), occurs in exactly row \(r_\ell\),
and every \(\ell>U\) occurs in no row. This proves the one-use incidence
claim for arbitrary features and labels, without any realizability premise.
\end{proof}

\begin{proposition}[Total executable improper one-block simulator]\label{prop:step-010-simulator}
Under Assumption~\ref{assump:candidate-regime}, the contradiction hypothesis~\eqref{eq:local-negation},
Propositions~\ref{prop:step-004-finite-hard-prior},
\ref{prop:step-005-certificate}, and
\ref{prop:step-009-overflow}, Lemma~\ref{lem:step-007-restriction-legality}, and
Lemma~\ref{lem:step-010-public-preprocessing},
Proposition~\ref{prop:step-010-row-construction}, and
Lemma~\ref{lem:step-010-one-use}, the following algorithm defines a
total randomized kernel
\[
  B_{\mu_{N,M},A}:\mathsf Z_M\longrightarrow\mathsf G_N
\]
for every arbitrary randomized \(A\) in (4).

On input \(z\in\mathsf Z_M\):

\begin{enumerate}
\item Sample the complete seed \(\Omega_{\mathrm{pre}}\) in the order stated
by Lemma~\ref{lem:step-010-public-preprocessing}.
\item Compute \(U\) from \(J,I_1,\ldots,I_n\), without evaluating \(z\).
\item If \(U>M\), return the fixed hypothesis
   \[
     g_0(x):=0\quad(x\in[N])
   \tag{16}
   \]
   immediately, without evaluating any coordinate of \(z\) and without
   invoking \(A\);
\item If \(U\le M\), construct the exact
   \(\mathcal P_{\Omega_{\mathrm{pre}}}(z)\in\mathsf S_n\) from (12), draw
   \[
     H\sim A(\mathcal P_{\Omega_{\mathrm{pre}}}(z))
   \]
   using \(A\)'s own fresh internal randomness, and return \(D_JH\).
\end{enumerate}

This simulator has exact input size \(M\), makes at most one call to \(A\)
and only with exact input size \(n\), uses each input row at most once, and
is defined on every arbitrary size-\(M\) labeled input. It does not know a
hidden \(t,Q\), does not require \(A\) to be proper or tag-symmetric, and
does not infer privacy or an ideal-law coupling.
\end{proposition}

\begin{proof}
The seed law is a valid input-independent probability kernel by
Lemma~\ref{lem:step-010-public-preprocessing}. On the overflow branch,
\(g_0\) is a fixed member of \(\mathsf G_N\); the control flow reaches (16)
before any input dereference or call to \(A\). Thus this branch is total and
input-independent even when the supplied input has arbitrary corrupt
labels.

On the nonoverflow branch,
Proposition~\ref{prop:step-010-row-construction} proves that the argument of
\(A\) is a valid member of its complete domain \(\mathsf S_n\) and has
exactly \(n\) rows. Therefore the randomized map \(A\) returns some
\(H\in\mathcal H_{k,N}\). Since \(J\in[k]\), Lemma~\ref{lem:step-007-restriction-legality} gives
\(D_JH\in\mathsf G_N\), regardless of whether \(H\) is a threshold,
monotone, proper, symmetric, deterministic, or tag-asymmetric.

Equivalently, if \(E\subseteq\mathsf G_N\) and
\(A(s,h)\) denotes the output probability assigned by \(A\) to
\(h\in\mathcal H_{k,N}\) on \(s\in\mathsf S_n\), then the simulator kernel
is the well-defined finite expectation

\[
B_{\mu_{N,M},A}(z,E)
=
\mathbb E_{\Omega_{\mathrm{pre}}}
\left[
\begin{cases}
  \mathbf1\{g_0\in E\}, & U>M,\\[1mm]
  \displaystyle
  \sum_{h\in\mathcal H_{k,N}}
    A(\mathcal P_{\Omega_{\mathrm{pre}}}(z),h)
    \mathbf1\{D_Jh\in E\}, & U\le M.
\end{cases}
\right].
\tag{17}
\]

Every term is nonnegative, the expression equals one when
\(E=\mathsf G_N\), and it is additive over disjoint events. Because all
input and output spaces are finite, (17) is a total randomized kernel on
every \(z\in\mathsf Z_M\). Formula (17) is a mathematical description of
the procedural branches; on \(U>M\) it contains no evaluation of
\(\mathcal P_\omega(z)\) and no \(A\)-term.

Lemma~\ref{lem:step-010-one-use} proves the exact access pattern. In
particular:

\begin{itemize}
\item At \(U=0\), no input row is read, all \(n\) rows come from the realized
  nonhidden instances, and \(A\) is called once on that size-\(n\) dataset;
\item At \(1\le U\le M\), precisely the first \(U\) input rows are read once
  each, in their original order, and all trailing rows are ignored;
\item At \(U=M\), all \(M\) rows are used once if that event is possible.
\item At \(U>M\), no row is read and \(A\) is not called.
\end{itemize}

The same clauses apply to endpoint labels, point-mass nonhidden
distributions, and arbitrary nonrealizable input labels. For \(k=2,3\),
Proposition~\ref{prop:step-009-overflow} says the overflow branch
has probability zero, but the branch remains defined. No step of this proof
uses Assumption~\ref{assump:central-dp}, compares two input datasets,
identifies an ideal sample law, or evaluates population risk. Those are
separate downstream obligations.
\end{proof}
